\noindent\textbf{Part 1.} \hfill(20 points)

\begin{itemize}
\item With the parallax we get the distance:
$D = \frac{1}{parallax} = 48.38$~pc, \hfill(1 point)\\
then we get the absolute magnitude using the distance modulus
$M_V = m_v - 5\log_{10}(D) + 5 = 4.2067$, \hfill(1 point)\\
and finally we go for the luminosity as BC is assumed to be the same for all
F/G stars, $(M_V - M_{V_\odot}) = (M_{bol} - M_{bol_\odot})$ \hfill(2 points)\\
by comparing with the magnitude of the Sun:
$\frac{L}{L_\odot} = 10^{-0.4(M_V - M_{V_\odot})} = 1.759$. \hfill(1 point)
\item We find the stellar radius from the Stephan--Boltzmann
% sic: source spelling "Stephan-Boltzmann"
law, which in solar units leads to
$\frac{T_{eff\,*}^4}{T_{eff\,\odot}^4}
 = \frac{R_\odot^2}{R_*^2}\frac{L_*}{L_\odot}$ \hfill(2 points)\\
i.e.,
$R_*[R_\odot] = \frac{T_{eff\,\odot}^2\sqrt{L_*\,[L_\odot]}}{T_{eff\,*}^2}
 = 1.238$. \hfill(2 points)
\item Now to find the stellar mass consider the surface gravitational
acceleration, which for Gauss' law and spherical symmetry considerations is:
$g = G\frac{M_*}{R_*^2}$, \hfill(1 point)\\
then one solves $M_*$ for and puts it in the right units,
$M_* = 1.944\ [M_\odot]$ \hfill(3 points)
\item Planet's orbital radius comes from
$a = \sqrt[3]{\frac{GMT^2}{4\pi^2}} = 0.0565$~A.U. \hfill(3 points)
\item Finally, checking the dimming due to the transit this is of
approximately 2.6\%, which is give
% sic: source wording "which is give by"
by the ratio between the areas of the planet disk and the stellar disk, i.e.,
$\left(\frac{R_p}{R_*}\right)^2 = 0.026$, \hfill(2 points)\\
leading to $R_p = 1.988\ [R_J]$ \hfill(2 points)
\end{itemize}

\begin{center}
\begin{tabular}{ccccc}
\toprule
Luminosity of & Radius of & Mass of & Mean planet's & Radius of the planet\\
the star & the star & the star & orbital radius & in Jupiter's radius\\
$L_*[L_\odot]$ & $R_*[R_\odot]$ & $M_*[M_\odot]$ & $a$[au] & $R_p[R_J]$\\
\midrule
1.759 & 1.24 & 1.94 & 0.0565 & 1.988\\
\bottomrule
\end{tabular}
\end{center}

Tolerance: of 5\% on the final answers is acceptable, numbers approximated to
the second digit in the partial calculations are OK.

% FIGURE NEEDED: fitted transit lightcurve -- normalized flux vs time with
% the trapezoidal model curve and red dashed tolerance envelope
% (2021_DA2_solution.pdf, page 1).

\medskip
\noindent\textbf{Part 2.} \hfill(25 points)

\noindent\textbf{(2.1)} \hfill(15 points)\\
One first needs to calibrate the axes. The shown circle representing the Earth
gives as a vertical mark in 5778~$^\circ$K,
% sic: source writes "5778 °K"
and a horizontal one at $S_{eff} = 1$. Then one need to retrieve additional
data from one or more of the additional observations:

\begin{itemize}
\item Having the temperature mark for the Sun, one can estimate $S_{min}$,
adding a second mark on the horizontal axis, so the scale of this one can be
calibrated. Then one can proceed backwards, starting from a given value of
$S_{min}$ or $S_{max}$, find the corresponding temperature, so the vertical
axis can be calibrated as well.
\item Alternatively, one can notice that the limiting $S_{max}$ curve crosses
the X axis in $S_{eff} = 1$, and look numerically for the corresponding
temperature of the bottom line, thus calibrating the vertical axis. Then
starting from a given temperature, one calibrates de
% sic: source wording
horizontal one pivoting on $S_{min}$ / $S_{max}$ calculations as explained
before.

There are several ways, some more precise than others. The student should be
as careful as possible (we give small tolerance to the error in the
calibrations of the axes). A tricky part is that the X-axis increases towards
the left, but this is something rigorous students will discover, as there is
no self-consistency otherwise. Intuition may also help if one knows that
Earth is closer to the high-flux end of the Sun's habitable zone.
\item The luminosity of Gorgona's host star was found in Part A. After
finding the luminosity of Amacayacu's one the final plot looks like this:
\end{itemize}

For placing Gorgona in the right position \hfill(2 pt)

% FIGURE NEEDED: habitable-zone diagram, T_eff (5400-6100 K) vs S_eff
% (X axis reversed, 2.0 -> 0.0) with the shaded HZ band between the dashed
% S_min and S_max curves and the calculated positions of Gorgona (square)
% and Amacayacu (star) marked (2021_DA2_solution.pdf, page 2).

For calculating flux and placing Amacayacu in the right position \hfill(4 pt)

Tolerance for the position of the planets: $\pm0.025$ in $S_{eff}$; $\pm25$
[K] in $T_{eff}$.

\medskip
\noindent\textbf{(2.2)} \hfill(10 points)\\
NONE of the exoplanets lie in habitable zones. They orbit their stars too
close, so the effective flux is many times larger than $S_{max}$. The
long-but-straightforward answer is to perform all the calculations. The
elegant, clever one, is to note that Amacayacu orbits the faintest/coolest
star, and it is the one farthest away, so its effective flux is the absolute
minimum of the whole set. This effective flux was calculated as ${\sim}0.3$
in A), assuming a distance of 1~A.U. With the real distance, 0.08~A.U, it
grows by a factor $\frac{1}{0.08^2}\sim156$, i.e., it has
$S_{eff}\sim46$, but the habitable zone for this range of temperatures can
not exceed $S_{eff}\sim1.1$, which completely closes the case.

\begin{center}
\begin{tabular}{lr}
\toprule
Name & $S_{eff}$\\
\midrule
Tayrona & 1785.994090\\
Iguaque & 1499.270574\\
Gorgona & 1096.175314\\
Amacayacu & 50.794890\\
Malpelo & 1374.231792\\
Pisba & 5568.747040\\
Tatama & 678.730709\\
\bottomrule
\end{tabular}
\end{center}

\medskip
\noindent\textbf{Part 3.}

\noindent\textbf{(3.1)} (3 points per each row; 1 point extra for the only
number in the last row)

\begin{center}
\begin{tabular}{cccccc}
\toprule
Low-mass sample & Sample size & $\mu$ & $\sigma$ & $\mu+\sigma$
  & No.\ of planets to exclude\\
\midrule
Full / Original & 38 & 1.99 & 2.55 & 4.54 & 4\\
1st subsample & 34 & 1.23 & 1.11 & 2.34 & 5\\
2nd subsample & 29 & 0.84 & 0.61 & 1.45 & 5\\
Final subsample & 24 & -- & -- & -- & --\\
\bottomrule
\end{tabular}
\end{center}

\noindent\textbf{(3.2)}
% FIGURE NEEDED: sorted planet masses -- Mass vs no. of planet in the list,
% blue dots increasing from 0 to ~12, with three horizontal red lines at the
% successive mu+sigma cuts (2021_DA2_solution.pdf, page 4).

\noindent\textbf{(3.3)} (Tolerance 5\%)

\begin{center}
\begin{tabular}{cccccc}
\toprule
$T_{eff}$ & Min. & 1st Quartile & Median & 3rd Quartile & Max\\
\midrule
LME & 5209 & 5574.5 & 5932.5 & 6181.75 & 6460\\
HME & 5548 & 5992.5 & 6177.0 & 6255.0 & 6490\\
\bottomrule
\end{tabular}
\end{center}

\noindent\textbf{(3.4)}
% FIGURE NEEDED: horizontal boxplots of T_eff for the HME (pink, with one
% outlier circle near 5550 K) and LME (green) samples, temperature axis
% 5200-6400 K (2021_DA2_solution.pdf, page 4).

Yes, the boxplots indicate that HMEs tend to get formed around hotter stars.
\hfill(1 point)
